%%%%%%%%%%%%%%%%%%%%%%%%%%%%%%%%%%%%%%%%%%%%%%
%%   华中师范大学硕士论文模板 - 宏包配置
%%   Package Configuration for CCNU Thesis
%%   版本: 2.0 (2025年3月更新)
%%
%%   文件用途说明:
%%   ============
%%   本文件集中配置论文所需的所有 LaTeX 宏包
%%   按功能分类组织，便于管理和维护
%%   
%%   注意: 宏包加载顺序很重要！
%%   一般原则: 基础宏包 -> 功能宏包 -> 依赖宏包
%%%%%%%%%%%%%%%%%%%%%%%%%%%%%%%%%%%%%%%%%%%%%%

%%===========================================================================
%%  中文支持宏包 (核心)
%%===========================================================================
%% ctex - 现代中文排版宏包
%% 用途: 提供完整的中文 LaTeX 支持
%% 为什么需要: 
%%   - 自动处理中文字体配置
%%   - 提供中文标题、章节编号
%%   - 支持中文标点符号
%% 关键选项说明:
%%   - UTF8: 使用 UTF-8 编码
%%   - heading=true: 启用章节标题样式
%%   - zihao=-4: 默认字号小四 (12pt)
%%   - fontset=fandol: 使用 Fandol 字体集
%%     (备选: windows, mac, ubuntu 等)
\usepackage[UTF8,heading=true,zihao=-4,fontset=fandol]{ctex}

%%===========================================================================
%%  页面布局宏包
%%===========================================================================
%% geometry - 页面尺寸和边距设置
%% 用途: 精确控制页面布局
%% 为什么需要: 学校对页边距有严格要求
%% 关键选项说明:
%%   - top/bottom/left/right: 四边边距
%%   - headsep: 页眉与正文距离
%%   - headheight: 页眉高度
%%   - footskip: 页脚与正文距离
\usepackage{geometry}

% titletoc: 目录格式控制（引导符、缩进等）
\usepackage{titletoc}

%%===========================================================================
%%  图形处理宏包
%%===========================================================================
%% graphicx - 插入图片
%% 用途: 支持各种格式图片插入
%% 为什么需要: 论文中需要插入图表
%% 关键选项说明:
%%   - 支持格式: PDF, PNG, JPG, EPS 等
%%   - 主要命令: \includegraphics[options]{file}
\usepackage{graphicx}

%% float - 浮动体控制
%% 用途: 控制图表浮动位置
%% 为什么需要: 精确控制图表位置，防止乱飘
%% 关键选项说明:
%%   - [H]: 强制放在当前位置
%%   - [htbp]: 自动选择最佳位置
\usepackage{float}

%%===========================================================================
%%  数学公式宏包
%%===========================================================================
%% amsmath - 美国数学学会数学排版
%% 用途: 提供高级数学排版功能
%% 为什么需要: 撰写理工科论文必备
%% 关键功能:
%%   - align: 对齐公式环境
%%   - equation*: 无编号公式
%%   - \text{}: 公式内文本
\usepackage{amsmath,amssymb,amsfonts}

%% bm - 粗体数学符号
%% 用途: 在数学环境中使用粗体
%% 为什么需要: 向量、矩阵常用粗体表示
%% 示例: \bm{v} 表示粗体向量 v
\usepackage{bm}

%% mathrsfs - 花体字母
%% 用途: 提供 \mathscr{} 花体命令
%% 为什么需要: 常用于表示变换、算子等
%% 示例: \mathscr{F} 表示傅里叶变换
\usepackage{mathrsfs}

%%===========================================================================
%%  表格处理宏包
%%===========================================================================
%% booktabs - 专业表格线
%% 用途: 提供 \toprule, \midrule, \bottomrule
%% 为什么需要: 比 \hline 更美观的表格线
%% 特点: 线条粗细自动调整
\usepackage{booktabs}

%% array - 表格列格式扩展
%% 用途: 扩展表格列对齐方式
%% 为什么需要: 支持自定义列格式
%% 示例: >{\centering\arraybackslash}p{3cm}
\usepackage{array}

%% multirow - 跨行单元格
%% 用途: 实现表格单元格跨多行
%% 为什么需要: 复杂表格布局
%% 示例: \multirow{2}{*}{内容}
\usepackage{multirow}

%% tabularx - 自动宽度表格
%% 用途: 表格宽度自动分配
%% 为什么需要: 需要精确控制表格宽度
%% 特点: X 列自动拉伸填充
\usepackage{tabularx}

%%===========================================================================
%%  图表标题宏包
%%===========================================================================
%% caption - 浮动体标题样式
%% 用途: 自定义图/表标题格式
%% 为什么需要: 学校对标题格式有要求
%% 关键选项:
%%   - font=small: 标题字号
%%   - labelsep=quad: 标签与标题间距
\usepackage{caption}

%% subcaption - 子图标题
%% 用途: 支持子图/子表编号
%% 为什么需要: 多子图并排显示
%% 示例: \begin{subfigure}{0.5\textwidth}
\usepackage{subcaption}

%%===========================================================================
%%  参考文献宏包
%%===========================================================================
%% natbib - 作者-年份和数字引用
%% 用途: 提供灵活的引用格式
%% 为什么需要: 支持 \citep, \citet 等命令
%% 关键选项:
%%   - numbers: 数字引用格式
%%   - sort&compress: 排序并压缩连续引用
%%   - square: 使用方括号
\usepackage[numbers,sort&compress,square]{natbib}

%%===========================================================================
%%  页眉页脚宏包
%%===========================================================================
%% fancyhdr - 自定义页眉页脚
%% 用途: 完全控制页眉页脚内容
%% 为什么需要: 学校要求特定格式
%% 关键命令:
%%   - \fancyhead: 设置页眉
%%   - \fancyfoot: 设置页脚
\usepackage{fancyhdr}

%%===========================================================================
%%  颜色宏包
%%===========================================================================
%% xcolor - 扩展颜色支持
%% 用途: 定义和使用颜色
%% 为什么需要: 彩色链接、高亮等
%% 预定义颜色: red, blue, green, black 等
\usepackage{xcolor}

%%===========================================================================
%%  超链接宏包 (最后加载)
%%===========================================================================
%% hyperref - 超链接和 PDF 元数据
%% 用途: 生成可点击链接和 PDF 书签
%% 为什么需要: 目录、引用可点击跳转
%% 注意: 应尽可能最后加载
%% 关键选项:
%%   - bookmarksnumbered: 书签带编号
%%   - colorlinks: 彩色链接 (打印时设为 false)
%%   - citecolor/linkcolor/urlcolor: 链接颜色
\usepackage[
    bookmarksnumbered=true,
    bookmarksopen=true,
    colorlinks=false,
    citecolor=blue,
    linkcolor=red,
    anchorcolor=green,
    urlcolor=blue,
    pdfencoding=auto
]{hyperref}

%%===========================================================================
%%  智能引用宏包 (hyperref 之后加载)
%%===========================================================================
%% cleveref - 智能引用
%% 用途: 自动添加 "图"、"表"、"公式" 前缀
%% 为什么需要: 简化引用，自动处理编号
%% 示例: \cref{fig:example} -> "图 1.1"
\usepackage{cleveref}

%%===========================================================================
%%  其他实用宏包
%%===========================================================================
%% indentfirst - 章节首段缩进
%% 用途: 章节第一段也缩进
%% 为什么需要: 中文排版习惯
\usepackage{indentfirst}

%% makeidx - 索引制作
%% 用途: 生成文档索引
%% 为什么需要: 长文档关键词索引
%% 使用: \makeindex 和 \printindex
\usepackage{makeidx}

%% footnote - 脚注增强
%% 用途: 改进脚注功能
%% 为什么需要: 更好的脚注控制
\usepackage{footnote}

%% footmisc - 脚注样式
%% 用途: 自定义脚注格式
%% 为什么需要: 每页重新编号脚注
%% 关键选项:
%%   - perpage: 每页重新编号
%%   - symbol: 使用符号而非数字
\usepackage[perpage,symbol]{footmisc}

%%===========================================================================
%%  数学排版优化
%%===========================================================================
%% 允许公式跨页断行
%% 用途: 避免长公式导致大量空白
%% 注意: 可能影响阅读连续性
\allowdisplaybreaks

%%===========================================================================
%%  文件结束
%%===========================================================================
